% !TEX program = xelatex
% =====================================================================
%  群体遗传学 · mtDNA 与 chrY —— 幻灯片内容（只改这个文件即可）
%  编译：xelatex main.tex（两遍，生成页码与进度条）
%  版式组件在 theme.tex；使用说明见 README.md
%  图片放在 image/ 目录，\PicLeft 里只写文件名，如 1.png
% =====================================================================
\documentclass[aspectratio=169]{ctexbeamer}
% =====================================================================
%  theme.tex —— NeonDark 深紫霓虹风主题（版式 + 组件库）
%  一般不需要改动；内容请只改 main.tex。
%  想换配色：只改下面「主题色板」里的 HTML 色值即可。
% =====================================================================
\usepackage{amsmath, amssymb}
\usepackage{tikz}
\usepackage{graphicx}
\graphicspath{{image/}}
\usepackage[most]{tcolorbox}
\usefonttheme[onlymath]{serif}
% ----------------------------- 主题色板 -----------------------------
\definecolor{deepbg}{HTML}{0A0618}        % 全局深底
\definecolor{bgcard}{HTML}{161029}    % 卡片底
\definecolor{bgdeep}{HTML}{0D0A1E}    % 更深的嵌套底
\definecolor{bgbadge}{HTML}{150E2C}   % 徽章底
\definecolor{border}{HTML}{2B2344}    % 卡片描边
\definecolor{lighttext}{HTML}{F0E6FF} % 高亮正文
\definecolor{bodytext}{HTML}{D4C8E8}  % 常规正文
\definecolor{muted}{HTML}{9A8AAA}     % 次要文字
\definecolor{fainttext}{HTML}{6A5A7E} % 弱化文字
\definecolor{purple}{HTML}{7C3AED}    % 主紫
\definecolor{violetL}{HTML}{A78BFA}   % 亮紫
\definecolor{blueL}{HTML}{60A5FA}     % 亮蓝
\definecolor{pink}{HTML}{EC4899}      % 粉
\definecolor{pinkL}{HTML}{F472B6}     % 亮粉
\definecolor{cyanL}{HTML}{22D3EE}     % 亮青
\definecolor{goldL}{HTML}{FBBF24}     % 亮金
\definecolor{redL}{HTML}{F87171}      % 亮红
\definecolor{greenL}{HTML}{34D399}    % 亮绿

% ----------------------------- 全局字体 -----------------------------
\setbeamerfont{normal text}{size=\small}
\setbeamerfont{frametitle}{size=\large, series=\bfseries}
\setbeamerfont{framesubtitle}{size=\footnotesize}

% ----------------------------- Beamer 基础设置 ----------------------
\setbeamercolor{background canvas}{bg=deepbg}
\setbeamercolor{normal text}{fg=bodytext}
\setbeamercolor{frametitle}{fg=violetL}
\setbeamercolor{framesubtitle}{fg=muted}
\setbeamercolor{math text displayed}{fg=lighttext}
\setbeamercolor{math text inline}{fg=lighttext}
\setbeamercolor{title}{fg=lighttext}
\setbeamertemplate{navigation symbols}{}
\setbeamertemplate{blocks}[rounded]
\setbeamertemplate{caption}[numbered]

% ----------------------------- 背景：暗紫光晕 + 星点 ----------------
% 模拟原 HTML 的 Canvas 粒子背景（静态、确定性）
\setbeamertemplate{background canvas}{%
  \begin{tikzpicture}[remember picture, overlay,
      shift={(current page.south west)}, x=\paperwidth, y=\paperheight]
    \fill[deepbg] (0,0) rectangle (1,1);
    % 角落光晕
    \fill[purple, opacity=0.06] (1.0,1.0) circle (0.38);
    \fill[pink,   opacity=0.05] (0.0,0.0) circle (0.32);
    \fill[blueL,  opacity=0.04] (1.0,0.05) circle (0.26);
    \fill[purple, opacity=0.05] (0.05,1.0) circle (0.22);
    % 星点
    \foreach \x/\y/\r in {0.10/0.86/2.2, 0.24/0.72/1.4, 0.40/0.93/1.8,
                          0.62/0.80/1.2, 0.78/0.90/2.0, 0.92/0.66/1.4,
                          0.86/0.30/1.8, 0.70/0.16/1.2, 0.55/0.06/1.6,
                          0.36/0.20/1.4, 0.18/0.10/1.8, 0.07/0.38/1.2,
                          0.94/0.12/1.2, 0.50/0.50/1.0}{
      \fill[violetL, opacity=0.10] (\x,\y) circle (\r pt);}
    \foreach \x/\y/\r in {0.16/0.62/1.2, 0.33/0.40/1.6, 0.68/0.58/1.4,
                          0.88/0.48/1.2, 0.45/0.80/1.0, 0.75/0.05/1.6,
                          0.05/0.55/1.4, 0.58/0.30/1.0}{
      \fill[blueL, opacity=0.08] (\x,\y) circle (\r pt);}
    \foreach \x/\y/\r in {0.28/0.88/1.0, 0.82/0.76/1.6, 0.65/0.35/1.2,
                          0.12/0.24/1.0, 0.40/0.60/1.0, 0.92/0.86/1.2}{
      \fill[pinkL, opacity=0.07] (\x,\y) circle (\r pt);}
  \end{tikzpicture}%
}

% ----------------------------- 标题区模板 ----------------------------
% 渐变分隔线：紫 → 粉
\newcommand{\gradrule}{%
  \tikz\fill[draw=none, left color=purple, right color=pink,
             rounded corners=1pt] (0,0) rectangle (1.35cm,0.09cm);}

\setbeamertemplate{frametitle}{%
  \vspace{1.6mm}%
  {\usebeamerfont{frametitle}\usebeamercolor[fg]{frametitle}\insertframetitle\par}%
  \vspace{1.0mm}\gradrule\vspace{0.6mm}%
  \ifx\insertframesubtitle\empty\else
    {\usebeamerfont{framesubtitle}\usebeamercolor[fg]{framesubtitle}\insertframesubtitle\par}%
  \fi
  \vspace{0.6mm}%
}

% ----------------------------- 页脚：页码 + 进度条 -------------------
\setbeamertemplate{footline}{%
  \begin{beamercolorbox}[ht=2.2mm, dp=1.0mm, center]{footline}
    \color{fainttext}\scriptsize
    \textcolor{violetL}{$\blacktriangleleft$}\hspace{1.5mm}%
    \insertframenumber\,/\,\inserttotalframenumber\hspace{1.5mm}%
    \textcolor{pinkL}{$\blacktriangleright$}
  \end{beamercolorbox}%
  \begin{tikzpicture}
    \fill[border] (0,0) rectangle (\paperwidth, 1.2pt);
    \fill[draw=none, left color=purple, right color=pink]
      (0,0) rectangle
      (\dimexpr\paperwidth*\insertframenumber/\inserttotalframenumber\relax, 1.2pt);
  \end{tikzpicture}%
}

% =====================================================================
%  模板组件库（自定义命令，可直接复用）
% =====================================================================
% \Badge{描边色}{文字色}{文字} —— 胶囊徽章
\newcommand{\Badge}[3]{%
  \tcbox[on line, colback=bgbadge, colframe=#1, boxrule=0.6pt, arc=8pt,
         left=7pt, right=7pt, top=2pt, bottom=2pt,
         before skip=2pt, after skip=2pt]
    {\color{#2}\footnotesize\bfseries #3}}

% 卡片环境：card（普通）/ cardc{描边色}（彩色描边卡片）
\newtcolorbox{card}{%
  colback=bgcard, colframe=border, boxrule=0.6pt, arc=3mm,
  left=3.5mm, right=3.5mm, top=2.2mm, bottom=2.2mm,
  before skip=1.6mm, after skip=1.6mm}
\newtcolorbox{cardc}[1]{%
  colback=bgcard, colframe=#1, boxrule=0.7pt, arc=3mm,
  left=3.5mm, right=3.5mm, top=2.2mm, bottom=2.2mm,
  before skip=1.6mm, after skip=1.6mm}

% 提示条：\notebox{描边色}{文字色}{内容}
\newcommand{\notebox}[3]{%
  \begin{tcolorbox}[colback=bgdeep, colframe=#1, boxrule=0.5pt, arc=2mm,
                    left=4mm, right=4mm, top=1.6mm, bottom=1.6mm,
                    before skip=1.6mm, after skip=1.6mm]
    \centering{\color{#2}\small #3}
  \end{tcolorbox}}

% 卡片标题与描述的简写
\newcommand{\cardtitle}[1]{{\color{lighttext}\bfseries #1}}
\newcommand{\carddesc}[1]{{\color{muted}\footnotesize #1}}

% 地图节点：\fnode{主文字}{副文字}（普通）/ \fnodeC{主文字}（核心，紫底高亮）
\newcommand{\fnode}[2]{%
  \tcbox[on line, colback=bgcard, colframe=border, boxrule=0.6pt, arc=2mm,
         left=5pt, right=5pt, top=2.5pt, bottom=2.5pt]
    {{\color{bodytext}#1}\ {\color{fainttext}\scriptsize #2}}}
\newcommand{\fnodeC}[1]{%
  \tcbox[on line, colback=purple!22!deepbg, colframe=purple, boxrule=0.7pt, arc=2mm,
         left=5pt, right=5pt, top=2.5pt, bottom=2.5pt]
    {\color{lighttext}\bfseries #1}}

% 流程步骤节点：\stepnode{章号}{名称}{描述} / \stepnodeH（高亮）
\newcommand{\stepnode}[3]{%
  \tcbox[on line, colback=bgcard, colframe=border, boxrule=0.6pt, arc=2mm,
         left=3pt, right=3pt, top=2pt, bottom=2pt]
    {\begin{minipage}{1.72cm}\centering
     {\color{violetL}\scriptsize\bfseries #1}\\
     {\color{lighttext}\small\bfseries #2}\\
     {\color{fainttext}\tiny #3}
     \end{minipage}}}
\newcommand{\stepnodeH}[3]{%
  \tcbox[on line, colback=purple!22!deepbg, colframe=purple, boxrule=0.7pt, arc=2mm,
         left=3pt, right=3pt, top=2pt, bottom=2pt]
    {\begin{minipage}{1.72cm}\centering
     {\color{violetL}\scriptsize\bfseries #1}\\
     {\color{lighttext}\small\bfseries #2}\\
     {\color{fainttext}\tiny #3}
     \end{minipage}}}



% =====================================================================
%  左图右文组件
% =====================================================================
% 图片卡片：\PicCard{图片文件}{图注}
\newcommand{\PicCard}[2]{%
  \begin{tcolorbox}[colback=bgcard, colframe=border, boxrule=0.6pt, arc=3mm,
                    boxsep=0pt, left=1.6mm, right=1.6mm, top=1.6mm, bottom=1.6mm,
                    before skip=0pt, after skip=1mm]
    \centering
    \includegraphics[width=\linewidth, height=0.40\textheight, keepaspectratio]{#1}\\[1mm]
    {\color{fainttext}\scriptsize #2}
  \end{tcolorbox}}

% 要点条目：\Point{强调色}{标题}{说明}
\newcommand{\Point}[3]{%
  \begin{tcolorbox}[colback=bgcard, colframe=border, boxrule=0pt,
                    borderline west={2.2pt}{0pt}{#1}, arc=1.5mm,
                    left=3.2mm, right=2mm, top=0.9mm, bottom=0.9mm,
                    before skip=0.8mm, after skip=0.8mm]
    {\color{#1}\bfseries\small #2}\\[0.4mm]
    {\color{muted}\footnotesize #3}
  \end{tcolorbox}}

% 左图右文整页：\PicLeft{图片}{图注}{右侧内容}（右侧内容通常是若干 \Point）
\newcommand{\PicLeft}[3]{%
  \begin{columns}[T]
    \begin{column}{0.44\textwidth}\PicCard{#1}{#2}\end{column}
    \begin{column}{0.53\textwidth}#3\end{column}
  \end{columns}}

% 右图左文版本（想换个节奏时使用，参数相同）
\newcommand{\PicRight}[3]{%
  \begin{columns}[T]
    \begin{column}{0.53\textwidth}#3\end{column}
    \begin{column}{0.44\textwidth}\PicCard{#1}{#2}\end{column}
  \end{columns}}


% ----------------------------- 基本信息 -----------------------------
\title{群体遗传学：mtDNA 与 chrY}
\author{报告人姓名}
\institute{单位 · 学院 / 实验室}
\date{2026 年 11 月}

\begin{document}

% ---------------------------------------------------------------------
%  1 · 封面
% ---------------------------------------------------------------------
\begin{frame}[plain]
  \centering
  \vspace{6mm}
  \Badge{purple}{violetL}{单位 · 学院 / 实验室}
  \par\vspace{6mm}
  {\Huge\bfseries {\color{pinkL}mtDNA}~{\color{violetL}$\times$}~{\color{blueL}chrY}}
  \par\vspace{3.5mm}
  {\large\color{bodytext}\ziju{0.6}群体遗传学}
  \par\vspace{2.5mm}
  {\color{violetL}母系与父系，两条读取人群历史的“单行道”}
  \par\vspace{2.5mm}
  \gradrule
  \par\vspace{3.5mm}
  {\color{muted}报告人姓名 \quad·\quad 2026 年 11 月}
  \par\vspace{7mm}
  {\color{muted}\footnotesize
    线粒体 DNA\quad$\bullet$\quad Y 染色体\quad$\bullet$\quad
    单倍群\quad$\bullet$\quad 人群迁徙}
\end{frame}

% ---------------------------------------------------------------------
%  2 · 引入问题（三张卡片）
% ---------------------------------------------------------------------
\begin{frame}{为什么盯着这两段 DNA？}{三个常常被问到的问题}
  \Badge{pink}{pinkL}{问题引入}\vspace{1mm}

  \begin{columns}[t]
    \begin{column}{0.32\textwidth}
      \begin{card}
        \centering
        {\color{pinkL}\Large 母}\\[1mm]
        \cardtitle{我的母系祖先从哪来？}\\[1mm]
        \carddesc{mtDNA 只经母亲传递，几乎不重组}
      \end{card}
    \end{column}
    \begin{column}{0.32\textwidth}
      \begin{card}
        \centering
        {\color{blueL}\Large 父}\\[1mm]
        \cardtitle{父系祖先走过哪些路？}\\[1mm]
        \carddesc{chrY 非重组区只经父亲传递}
      \end{card}
    \end{column}
    \begin{column}{0.32\textwidth}
      \begin{card}
        \centering
        {\color{violetL}\Large $\Delta t$}\\[1mm]
        \cardtitle{这段历史有多古老？}\\[1mm]
        \carddesc{突变积累 $\to$ 分子钟 $\to$ 分化时间}
      \end{card}
    \end{column}
  \end{columns}

  \notebox{border}{pinkL}{\small
    单一世系的传递路径清晰 —— 数“突变”就能重建“谱系”}
\end{frame}

% ---------------------------------------------------------------------
%  3 · 对比（左右两张卡片）
% ---------------------------------------------------------------------
\begin{frame}{两条世系 · 一张对照表}{同样是“单倍体”，脾气却不一样}
  \Badge{purple}{violetL}{核心对比}\vspace{0.5mm}

  \begin{columns}[t]
    \begin{column}{0.49\textwidth}
      \begin{cardc}{pink!60!border}
        \centering
        {\color{pinkL}\Large mtDNA}\\[1mm]
        {\color{pinkL}\bfseries 母系 · 线粒体}\\[1.5mm]
        {\color{muted}\footnotesize
          环状，约 16.6 kb，37 个基因\\
          高拷贝、高突变率、可见异质性\\[1mm]
          {\color{pinkL}$\to$ 适合：深度祖源、古 DNA}}
      \end{cardc}
    \end{column}
    \begin{column}{0.49\textwidth}
      \begin{cardc}{blueL!60!border}
        \centering
        {\color{blueL}\Large chrY}\\[1mm]
        {\color{blueL}\bfseries 父系 · Y 染色体}\\[1.5mm]
        {\color{muted}\footnotesize
          线状，约 57 Mb，非重组区（NRY）占大部分\\
          拷贝少、突变率低、SNP + STR 双标记\\[1mm]
          {\color{blueL}$\to$ 适合：父系扩张、姓氏与家族}}
      \end{cardc}
    \end{column}
  \end{columns}

  \notebox{purple!50!border}{violetL}{\small
    两者合看：{\color{lighttext}谁迁徙了？男性还是女性？} \quad
    {\color{fainttext}\footnotesize $\to$ 性别偏向的人口历史}}
\end{frame}

% ---------------------------------------------------------------------
%  4–10 · 左图右文（把图片换成自己的图，把文字换成自己的解读）
% ---------------------------------------------------------------------
\begin{frame}{mtDNA：母系的“活化石”}{结构与遗传特点}
  \Badge{pink}{pinkL}{母系遗传}\vspace{0.5mm}

  \PicLeft{1.png}{图 1 · 占位图：mtDNA 环状图谱（替换为自己的图）}{
    \Point{pinkL}{只经母亲传递}{精子线粒体在受精后基本被清除，世系单线延续}
    \Point{violetL}{几乎无重组}{整条分子作为一个单倍型传递，系统树可直接解读}
    \Point{blueL}{高拷贝 · 高突变率}{每个细胞数百至数千拷贝，古 DNA 中更易保存}
    \Point{goldL}{异质性（heteroplasmy）}{同一个体可混有多种 mtDNA 序列，需要注意}
  }
\end{frame}

\begin{frame}{chrY：父系的“姓氏”}{结构与遗传特点}
  \Badge{blueL}{blueL}{父系遗传}\vspace{0.5mm}

  \PicLeft{2.png}{图 2 · 占位图：Y 染色体结构示意（替换为自己的图）}{
    \Point{blueL}{只经父亲传递}{非重组区（NRY）整体作为单倍型，从父亲传给儿子}
    \Point{violetL}{SNP 定单倍群，STR 看近亲}{SNP 稳定，适合划分深层分支；STR 变化快，适合家族尺度}
    \Point{cyanL}{高度重复的区域}{回文与扩增子序列使比对和分型更具挑战}
    \Point{pinkL}{与姓氏高度相关}{父系传递的模式与许多文化中姓氏的传承相似}
  }
\end{frame}

\begin{frame}{单倍群与系统发育树}{用突变为世系“命名”}
  \Badge{purple}{violetL}{世系分类}\vspace{0.5mm}

  \PicLeft{4.jpg}{图 3 · 占位图：单倍群系统树（替换为自己的图）}{
    \Point{violetL}{单倍群 = 共享突变的分支}{同一分支的个体拥有共同祖先和相同的标志性突变}
    \Point{pinkL}{mtDNA：L0–L6, M, N, R …}{以 rCRS 为参照，PhyloTree 维护标准命名}
    \Point{blueL}{chrY：A–T 与亚支}{ISOGG 维护命名，新 SNP 不断细化分支}
    \Point{goldL}{根在非洲}{两种世系最深的分支（L0 / A00, A0）均出现在非洲人群}
  }
\end{frame}

\begin{frame}{突变率与分子钟}{把“突变数”换算成“年代”}
  \Badge{goldL}{goldL}{时间标尺}\vspace{0.5mm}

  \PicLeft{5.png}{图 4 · 占位图：分子钟 / 分化时间估计（替换为自己的图）}{
    \Point{goldL}{基本思想}{突变大致随时间线性积累：\ $t \approx d\,/\,(2\mu)$}
    \Point{pinkL}{mtDNA 速率}{约 $10^{-8}$ /位点/年量级，控制区高于编码区}
    \Point{blueL}{chrY 速率}{约 $10^{-9}$ /位点/年量级，可用古 DNA 校准}
    \Point{redL}{注意校准}{速率不同，年代可差数倍，须注明所用速率}
  }
\end{frame}

\begin{frame}{有效群体大小与遗传漂变}{为什么这两段 DNA “更敏感”}
  \Badge{cyanL}{cyanL}{漂变更强}\vspace{0.5mm}

  \PicLeft{6.png}{图 5 · 占位图：瓶颈 / 奠基者效应示意（替换为自己的图）}{
    \Point{cyanL}{有效群体大小仅为常染色体的 1/4}{单倍体、单性别传递 $\to$ $N_e$ 更小}
    \Point{violetL}{漂变更快}{瓶颈、奠基者效应会在单一世系上留下清晰印记}
    \Point{pinkL}{扩张的“星状”信号}{短期内的快速扩张，在网络图中呈现星状辐射}
    \Point{goldL}{代价}{只反映一条世系，不代表全基因组历史}
  }
\end{frame}

\begin{frame}{古 DNA 与迁徙路线}{把现代世系放回历史坐标}
  \Badge{greenL}{greenL}{人群迁徙}\vspace{0.5mm}

  \PicLeft{1.png}{图 6 · 占位图：迁徙路线与古 DNA 样本（替换为自己的图）}{
    \Point{greenL}{现代 + 古代一起看}{古 DNA 提供时间切片，现代数据提供分布全景}
    \Point{blueL}{单倍群频率的时空变化}{某些单倍群的兴衰对应农业扩张、草原迁徙等事件}
    \Point{pinkL}{TMRCA（最近共同祖先）}{“线粒体夏娃”与“Y 染色体亚当”不是同一时代、同一个人}
    \Point{goldL}{结合考古与语言学}{单一证据很难定案，多源互证更可靠}
  }
\end{frame}

\begin{frame}{性别偏向的人口历史}{mtDNA 与 chrY 的“不同步”}
  \Badge{pink}{pinkL}{关键洞察}\vspace{0.5mm}

  \PicLeft{2.png}{图 7 · 占位图：mtDNA 与 chrY 多样性对比（替换为自己的图）}{
    \Point{pinkL}{父居 vs 母居}{从夫居人群的 chrY 地理结构常强于 mtDNA}
    \Point{blueL}{男性有效群体大小}{少数男性贡献大量后代，chrY 多样性骤降}
    \Point{violetL}{混血的不对称}{殖民、征服使两条世系的外来成分比例不同}
    \Point{goldL}{解读方法}{对比 $F_{ST}$、$\pi$ 与世系组成，再结合历史背景}
  }
\end{frame}

\begin{frame}{常用统计量}{描述多样性与分化}
  \Badge{violetL}{violetL}{统计指标}\vspace{0.5mm}

  \PicRight{4.jpg}{图 8 · 占位图：多样性 / 分化结果示意（替换为自己的图）}{
    \Point{violetL}{单倍型多样性 $Hd$ 与核苷酸多样性 $\pi$}{衡量群体内的变异水平}
    \Point{cyanL}{Tajima's $D$、Fu's $F_S$}{检验中性与群体扩张}
    \Point{pinkL}{$F_{ST}$ 与 AMOVA}{衡量群体间的分化，及变异的分层来源}
    \Point{goldL}{PCA / MDS / 网络图}{直观展示群体之间的亲缘关系}
  }
\end{frame}

% ---------------------------------------------------------------------
%  分析流程（流程节点）
% ---------------------------------------------------------------------
\begin{frame}{典型分析流程}{从原始数据到可解读的结论}
  \Badge{purple}{violetL}{分析主线}\vspace{0.5mm}

  \begin{center}
    \stepnodeH{步骤 1}{数据}{WGS / 芯片 / 捕获}\hspace{1pt}
    \stepnode{步骤 2}{质控}{比对 · 覆盖度 · 污染}\hspace{1pt}
    \stepnode{步骤 3}{分型}{单倍群 · 变异检出}\hspace{1pt}
    \stepnode{步骤 4}{建树}{ML · 贝叶斯 · 网络}\hspace{1pt}
    \stepnode{步骤 5}{统计}{多样性 · 分化 · 中性检验}\hspace{1pt}
    \stepnode{步骤 6}{解读}{年代 · 迁徙 · 历史}
  \end{center}

  \begin{columns}[t]
    \begin{column}{0.32\textwidth}
      \begin{tcolorbox}[colback=purple!14!deepbg, colframe=border, boxrule=0.5pt, arc=2mm,
                        left=2mm, right=2mm, top=1.5mm, bottom=1.5mm]
        \centering
        {\color{violetL}\scriptsize\bfseries 前两步}\\
        {\color{bodytext}\footnotesize 数据质量决定上限}\\
        {\color{fainttext}\tiny 古 DNA 尤其要看污染与损伤}
      \end{tcolorbox}
    \end{column}
    \begin{column}{0.32\textwidth}
      \begin{tcolorbox}[colback=pink!14!deepbg, colframe=border, boxrule=0.5pt, arc=2mm,
                        left=2mm, right=2mm, top=1.5mm, bottom=1.5mm]
        \centering
        {\color{pinkL}\scriptsize\bfseries 中间两步}\\
        {\color{bodytext}\footnotesize 分型与建树是核心}\\
        {\color{fainttext}\tiny 结果须与参考树交叉验证}
      \end{tcolorbox}
    \end{column}
    \begin{column}{0.32\textwidth}
      \begin{tcolorbox}[colback=cyanL!14!deepbg, colframe=border, boxrule=0.5pt, arc=2mm,
                        left=2mm, right=2mm, top=1.5mm, bottom=1.5mm]
        \centering
        {\color{cyanL}\scriptsize\bfseries 最后两步}\\
        {\color{bodytext}\footnotesize 统计与解读要克制}\\
        {\color{fainttext}\tiny 单一世系 $\neq$ 全部历史}
      \end{tcolorbox}
    \end{column}
  \end{columns}
\end{frame}

% ---------------------------------------------------------------------
%  工具与数据库（六宫格卡片）
% ---------------------------------------------------------------------
\begin{frame}[shrink=6]{常用工具与数据库}{替换成自己实际用到的软件}
  \Badge{cyanL}{cyanL}{工具箱}\vspace{0.5mm}

  \begin{columns}[t]
    \begin{column}{0.32\textwidth}
      \begin{card}
        \centering
        {\color{pinkL}\bfseries mtDNA}\\[1mm]
        \cardtitle{HaploGrep · PhyloTree}\\[0.8mm]
        \carddesc{mtDNA 单倍群分型与标准参考树}
      \end{card}
    \end{column}
    \begin{column}{0.32\textwidth}
      \begin{card}
        \centering
        {\color{blueL}\bfseries chrY}\\[1mm]
        \cardtitle{Yleaf · yhaplo · ISOGG}\\[0.8mm]
        \carddesc{Y 单倍群分型与命名体系}
      \end{card}
    \end{column}
    \begin{column}{0.32\textwidth}
      \begin{card}
        \centering
        {\color{goldL}\bfseries Tree}\\[1mm]
        \cardtitle{IQ-TREE · BEAST}\\[0.8mm]
        \carddesc{最大似然建树与贝叶斯分化时间估计}
      \end{card}
    \end{column}
  \end{columns}
  \begin{columns}[t]
    \begin{column}{0.32\textwidth}
      \begin{card}
        \centering
        {\color{violetL}\bfseries Net}\\[1mm]
        \cardtitle{PopART · Network}\\[0.8mm]
        \carddesc{中接网络、星状扩张可视化}
      \end{card}
    \end{column}
    \begin{column}{0.32\textwidth}
      \begin{card}
        \centering
        {\color{cyanL}\bfseries Stat}\\[1mm]
        \cardtitle{Arlequin · DnaSP}\\[0.8mm]
        \carddesc{多样性、Tajima's $D$、AMOVA、$F_{ST}$}
      \end{card}
    \end{column}
    \begin{column}{0.32\textwidth}
      \begin{card}
        \centering
        {\color{greenL}\bfseries Data}\\[1mm]
        \cardtitle{AADR · 1000 Genomes}\\[0.8mm]
        \carddesc{古代与现代人群的公开数据集}
      \end{card}
    \end{column}
  \end{columns}

  \centering{\color{fainttext}\scriptsize
    工具版本与参数请在正文或附录中写明，便于复现}
\end{frame}

% ---------------------------------------------------------------------
%  局限与避坑（三栏 + 两个提示框）
% ---------------------------------------------------------------------
\begin{frame}[shrink=8]{局限与避坑}{单一世系的力量，也是它的边界}
  \Badge{redL}{redL}{谨慎解读}\vspace{0.5mm}

  \begin{columns}[t]
    \begin{column}{0.32\textwidth}
      \begin{card}
        \centering
        {\color{violetL}\large\bfseries 1}\\[1mm]
        \cardtitle{只是一条世系}\\[0.8mm]
        \carddesc{mtDNA 与 chrY 各自只占祖先的一小部分\\
                  {\color{fainttext}需要常染色体数据互证}}
      \end{card}
    \end{column}
    \begin{column}{0.32\textwidth}
      \begin{card}
        \centering
        {\color{pinkL}\large\bfseries 2}\\[1mm]
        \cardtitle{样本与采样偏差}\\[0.8mm]
        \carddesc{样本量少、地域不均会扭曲频率\\
                  {\color{fainttext}注意亲缘个体与代表性}}
      \end{card}
    \end{column}
    \begin{column}{0.32\textwidth}
      \begin{card}
        \centering
        {\color{blueL}\large\bfseries 3}\\[1mm]
        \cardtitle{时间估计不确定}\\[0.8mm]
        \carddesc{突变率、世代时间都有分歧\\
                  {\color{fainttext}报告置信区间而非单一数字}}
      \end{card}
    \end{column}
  \end{columns}

  \begin{columns}[t]
    \begin{column}{0.49\textwidth}
      \begin{tcolorbox}[colback=redL!10!deepbg, colframe=redL, boxrule=0pt,
                        borderline west={2.5pt}{0pt}{redL}, arc=1.5mm,
                        left=3mm, right=3mm, top=1.5mm, bottom=1.5mm]
        {\color{redL}\bfseries\small 常见误区}\\[0.5mm]
        {\color{muted}\footnotesize
          “单倍群 = 民族 / 血统” $\to$ 单倍群只描述一条世系，并不等于个人的全部祖源}
      \end{tcolorbox}
    \end{column}
    \begin{column}{0.49\textwidth}
      \begin{tcolorbox}[colback=blueL!10!deepbg, colframe=blueL, boxrule=0pt,
                        borderline west={2.5pt}{0pt}{blueL}, arc=1.5mm,
                        left=3mm, right=3mm, top=1.5mm, bottom=1.5mm]
        {\color{blueL}\bfseries\small 最佳实践}\\[0.5mm]
        {\color{muted}\footnotesize
          mtDNA + chrY + 常染色体联合分析，并与考古、语言学证据相互印证}
      \end{tcolorbox}
    \end{column}
  \end{columns}

  \notebox{purple!50!border}{violetL}{\small
    让数据说话，但不要让它说得比它能说的更多}
\end{frame}

% ---------------------------------------------------------------------
%  结语
% ---------------------------------------------------------------------
\begin{frame}[plain]
  \centering
  \vfill
  {\color{goldL}\Large $\star$}
  \par\vspace{3.5mm}
  {\large\color{bodytext}母系与父系，两条世系，}
  \par\vspace{2mm}
  {\large 共同讲述{\color{lighttext}\bfseries 人类走过的路}。}
  \par\vspace{3.5mm}
  \gradrule
  \par\vspace{3.5mm}
  {\color{violetL}\ziju{0.2}单位 · 学院 / 实验室}
  \par\vspace{2mm}
  {\color{muted}\small 谢谢大家，欢迎提问与讨论}
  \par\vspace{5mm}
  {\color{fainttext}\footnotesize
    mtDNA\quad$\bullet$\quad chrY\quad$\bullet$\quad 群体遗传学\quad$\bullet$\quad
    2026}
  \vfill
\end{frame}

\end{document}
